%% Anotações da lição 01 — exemplo de como usar o arquivo de notas
\begin{campo}{Ideia central}
(Escreva aqui, com as suas palavras, a resposta da lição.)
\end{campo}

\begin{campo}{Texto em foco: 2 Timóteo 3:16–17}
A lição cita \rb{2Tm 3:16}. O contexto (\rb{3:14–17}) é o conselho de Paulo a
Timóteo, que conhecia os \enquote{escritos sagrados} desde a infância
(\rb{3:15}).

\interlinear[modo=leitura, titulo={2Tm 3:16–17 · SBLGNT}, legenda=false]{dados/biblia/2tm03_16-17.grc.tsv}

\begin{lexico}{θεόπνευστος}{G2315}{1}
Adjetivo formado por \gr{θεός} (Deus) e o verbo \gr{πνέω} (soprar):
\enquote{soprado por Deus}. Só aparece aqui no NT. O versículo lista quatro
finalidades com \gr{πρός}: ensino, repreensão, correção e disciplina na
justiça. \gr{ἐπανόρθωσις} (correção) e \gr{ἄρτιος} (completo, apto) também
ocorrem uma única vez no NT.
\end{lexico}
\end{campo}

\EBlivcampo[3]{Dúvidas para pesquisar}{}
\EBlivcampo[3]{Aplicação}{}
